%% =====================================================================
%%  T-04-01  The Thumbscrew
%%  -------------------------------------------------------------------
%%  One idea per file. Core is mandatory. Slots in canonical order.
%%  Max 700 words. No layout commands. No filler. New source material
%%  for this entry goes in sources/<BOOK>/T-04-01.tex -- never by
%%  rewriting this file (docs/RULES.md R0.1).
%% =====================================================================
%% @kind: topic
%% @id: T-04-01
%% @title: The Thumbscrew
%% @subtitle: Everyone has one load-bearing insecurity
%% @chapter: c04
%% @order: 10
%% @status: stable
%% @sources: B3
%% @updated: 2026-09-19
%% @allow-rewrite: slot migration to the seven-question format (docs/RULES.md section 2)

\topic{T-04-01}{The Thumbscrew}{Everyone has one load-bearing insecurity}


\begin{core}
Every person has a single point --- usually an insecurity, an unmanaged emotion, a need, or a small private pleasure --- where ordinary resistance does not apply. Finding it is the whole of the preparatory work; once found, it can be loaded repeatedly and cheaply.
\end{core}

\begin{mechanism}
Resistance is distributed unevenly. Most of a person's defences are spread across the areas they have thought about, and concentrated where they have been hurt before. The unguarded point is unguarded precisely because it is not experienced as a weakness; it is experienced as a preference, a taste, a reasonable sensitivity. That is why it is invisible to its owner and obvious to an observer with no stake. Once identified it does not need to be attacked --- it only needs to be referred to. A person defending a point they do not know they are defending will concede almost anything elsewhere to keep it unexamined.
\end{mechanism}

\begin{application}
  \item Collect over time rather than probe. The point reveals itself in repetition, not in a single conversation.
  \item Test lightly --- one remark, then observe rather than follow up.
  \item Note what they defend when nothing is under attack.
  \item Once identified, refer to it rather than pressing it; the reference does the work.
\end{application}

\begin{feedback}
  \item One subject changes their register; everything else stays level.
  \item They mention the same person, achievement, or grievance far more often than its importance warrants.
  \item A small remark lands out of all proportion, or is deflected with unusual speed.
  \item They volunteer information about it unprompted, then steer away.
  \item Praise in one specific direction reliably produces compliance in another.
\end{feedback}

\begin{failure}
Finding the point takes time and repeated contact, so the technique is weak in one-off exchanges, which is most exchanges. It is also highly fallible: repetition and disproportion have ordinary causes --- grief, boredom, a genuine enthusiasm --- and misreading them produces confident errors. Used offensively it is the most damaging item in this chapter and the one most likely to produce the long-memory counterpart described in T-03-01.
\end{failure}

\begin{countermeasures}
  \item Audit yourself first. Ask which subject changes your register, and what you mention too often. That list is your own exposure.
  \item Flatten your reactions deliberately in the areas you know are sensitive, especially with someone new.
  \item Notice disproportionate attention to one part of your life from someone with no standing to give it. Interest that specific, that early, is reconnaissance.
  \item Give the point air. An insecurity that has been named aloud to a trusted person loses most of its leverage, because it can no longer be threatened silently.
\end{countermeasures}

\sources{B3}
\seesources
\seealso{T-03-01, T-04-03, T-04-04, T-22-01}
